\section*{第 \ref{ch:g1:healthy-habits} 章——健康习惯}

\begin{solution}{exo:g1:healthy-habits:1}
\omterm{def:g1:healthy-habits:germs}{病菌}是\omterm{def:g1:living-or-not:living}{生物}，小得肉眼看不见，到处都有；
有些会在进入身体后让我们生病。因为看不见，
手看上去干干净净，
其实仍可能带着它们。
\end{solution}

\begin{solution}{exo:g1:healthy-habits:2}
吃东西前、上完厕所后，
以及玩完户外游戏或接触\omterm{def:g1:animals-around-us:animal}{动物}后。
\end{solution}

\begin{solution}{exo:g1:healthy-habits:3}
慢慢数到 $20$（大约唱一首短歌的时间）。容易
忘记的地方是手指之间、拇指周围，
以及手背。
\end{solution}

\begin{solution}{exo:g1:healthy-habits:4}
吃完东西后，牙齿上会残留一些食物碎屑。口腔里的\omterm{def:g1:healthy-habits:germs}{病菌}
以这些残渣为食，尤其是甜的残渣，
慢慢地就会在牙齿上蛀出洞来，使牙齿疼痛。刷牙
能在“大餐”开始前把残渣清走。
\end{solution}

\begin{solution}{exo:g1:healthy-habits:5}
每天早饭后和睡觉前各刷一次，
每次大约和唱一首歌一样长，每颗牙齿的每一面都要刷到。
\end{solution}

\begin{solution}{exo:g1:healthy-habits:6}
\omterm{prop:g1:healthy-habits:needs}{多样的食物}（早餐吃水果、喝水）、运动（
户外玩耍、骑自行车）和\omterm{prop:g1:healthy-habits:needs}{睡眠}（夜里睡得久，早点
上床）。例子不限。
\end{solution}

\begin{solution}{exo:g1:healthy-habits:7}
\omterm{prop:g1:healthy-habits:needs}{睡眠}让身体休息、自我修复并生长。儿童需要更长的夜间\omterm{prop:g1:healthy-habits:needs}{睡眠}，
——远比成年人需要的长。
\end{solution}

\begin{solution}{exo:g1:healthy-habits:8}
例如：打哈欠、眼睛刺痛或发沉、无缘无故地
闹脾气。
\end{solution}

\begin{solution}{exo:g1:healthy-habits:9}
乳牙要用很多年才会脱落，在此之前要咀嚼食物；
乳牙上的蛀洞和成年人牙齿上的蛀洞一样疼；
而且下面等待长出的恒牙，需要一个干净、健康的口腔。
写出其中两个理由即可。
\end{solution}

\begin{solution}{exo:g1:healthy-habits:10}
闪粉代表\omterm{def:g1:healthy-habits:germs}{病菌}。握手、开门和共用铅笔，
会让\omterm{def:g1:healthy-habits:germs}{病菌}在手与手之间无形地传递；
而用肥皂洗手并慢慢数到 $20$，
就能把它们洗掉，闪粉和\omterm{def:g1:healthy-habits:germs}{病菌}一并去除。
\end{solution}
